\section{Audit target and sampling design}\label{sec:audit-target}

We audit a frozen binary classifier, including its decision threshold, on a
prespecified finite family of $G$ mutually exclusive intersectional cells.
Write $E=\mathbf{1}\{f(X)=0\}$ for a classification error among observations
with true label $Y=1$. The classifier's false-negative rate (FNR) in cell $g$
is $p_g=\Pr(E=1\mid Y=1,g)$. Its effective sample size $n_g$ counts
positive-label observations, rather than all cell members. Our target is a
one-sided difference relative to the count-weighted complement, with a fixed
practical tolerance $0<\Delta<1$.

In the fixed-stratum design, positive counts $n_1,\ldots,n_G$ are fixed and
errors are independent Bernoulli observations with constant rate $p_g$
within each cell. Let $N=\sum_h n_h$, $\pi_g=n_g/N$, and assume
$1-\pi_g>0$ so that a reference exists. Define
\begin{equation}
 \begin{aligned}
 q_g&=\sum_{h\ne g}\frac{\pi_h}{1-\pi_g}p_h,\\
 \delta_g&=p_g-q_g,\qquad H_{0g}:\delta_g\le\Delta.
 \end{aligned}
 \label{eq:theory-target}
\end{equation}
Although the target error count is binomial, the complement count generally
has a Poisson-binomial law. With $\bar p=\sum_h\pi_hp_h$, the equivalent
identity $\delta_g=(p_g-\bar p)/(1-\pi_g)$ also shows that an injected
increase in $p_g$ need not equal the resulting disparity. We use the actual
$\delta_g$ to identify harmed cells.

In an independent and identically distributed (IID) population audit, the same
comparison has a different interpretation: $q_g=\Pr(E=1\mid Y=1,g^c)$. Conditional on the two positive
counts for a fixed cell, the target and complement counts are independent
binomials. Our fixed-count simulations instead use
Eq.~\eqref{eq:theory-target}; a binomial reference assumption is unnecessary
for their confirmation test. Neither model automatically describes
dependent patients, households, or complex survey sampling.

Let $\mathcal H=\{g:\delta_g>\Delta\}$. For audit flags $F_g\in\{0,1\}$,
the familywise error rate (FWER) is
$\Pr\{\exists g\notin\mathcal H:F_g=1\}$. The audit missed-harm rate is
\begin{equation}
 \mathrm{MHR}=\frac{1}{|\mathcal H|}
      \sum_{g\in\mathcal H}\Pr(F_g=0), \qquad |\mathcal H|>0.
 \label{eq:theory-mhr}
\end{equation}
Thus, MHR concerns audit non-detection, whereas classifier FNR concerns
prediction errors. MHR is undefined when no cell is harmed. On a fully
supported design, estimation is assessed separately through macro risk
$G^{-1}\sum_g\mathbb E(\widehat p_g-p_g)^2$ and prevalence-weighted risk
$\sum_g\pi_g\mathbb E(\widehat p_g-p_g)^2$. These averages do not specify
the probability of identifying a particular harmed cell.

Splitting must preserve the target. We assign independent error
observations to discovery $D$ and confirmation $C$, with
$n_h^C=c n_h$ exactly in every supported cell, $0<c<1$, and integer counts in both
parts. Then complement proportions, and hence $q_g$, are unchanged.
Random within-cell split counts generally change the conditional
complement weights. Such a design requires an explicitly standardized
reference instead. A zero target or reference count is recorded as
insufficient information, with a confirmation p-value of one; the cell remains
in the registered family. Non-rejection is inconclusive and does not
certify fairness.

\section{Estimation and screening can separate}\label{sec:separation}

Sharing information across cells can improve mean estimation, as pursued
by structured regression and SureMap~\citep{herlihy2024structured,khodak2024suremap}. To isolate the decision
mechanism, we first examine a fixed Beta posterior mean. This calculation
does not characterize all fitted empirical Bayes, structured regression,
or SureMap estimators. Its reference is known, and its threshold decision
is a descriptive screen rather than an error-controlled test.

\begin{proposition}[Fixed-prior risk and strict-screen boundary]
\label{prop:theory-shrinkage}
Let $X\sim\mathrm{Bin}(n,p)$, $n\ge1$, and fix $\kappa>0$ and
$\mu\in(0,1)$ independently of $X$. Set
$\widehat p=X/n$ and $\widetilde p=(X+\kappa\mu)/(n+\kappa)$.
For an oracle reference $r$, let $t=r+\Delta\in(0,1)$ and define
$A_0=\{\widehat p>t\}$ and $A_\kappa=\{\widetilde p>t\}$.
Writing $R_\kappa(p)=\mathbb E(\widetilde p-p)^2$ and
$R_0(p)=p(1-p)/n$, we have
\begin{equation}
 \begin{aligned}
 &R_\kappa(p)-R_0(p)\\
 &\quad=\frac{\kappa^2(p-\mu)^2-(2\kappa+\kappa^2/n)p(1-p)}
 {(n+\kappa)^2}.
 \end{aligned}
 \label{eq:theory-risk}
\end{equation}
Moreover,
\begin{equation}
 \begin{aligned}
 \Pr_p(A_\kappa)&=\Pr_p(X\ge K_\kappa),\\
 K_\kappa&=\lfloor nt+\kappa(t-\mu)\rfloor+1,
 \end{aligned}
 \label{eq:theory-cutoff}
\end{equation}
where $K_\kappa\le0$ or $K_\kappa>n$ gives probability one or zero,
respectively. If $t>\mu$, the integer cutoff cannot decrease. If
\begin{equation}
 \kappa(t-\mu)\ge n(1-t),
 \label{eq:theory-collapse}
\end{equation}
the shrunk screen never selects the cell. If $p=t+\varepsilon\le1$,
$\varepsilon>0$, and $a=\kappa(t-\mu)-n\varepsilon>0$, then
\begin{equation}
 \begin{aligned}
 \Pr_p(A_\kappa)&\le e^{-2a^2/n},\\
 \Pr_p(A_0^c)&\le e^{-2n\varepsilon^2}.
 \end{aligned}
 \label{eq:theory-screen-bound}
\end{equation}
\end{proposition}

Equation~\eqref{eq:theory-risk} balances variance reduction against bias;
Eq.~\eqref{eq:theory-cutoff} tracks the error count needed to cross a
decision threshold. Their joint dependence on effective count, tolerance,
signal, and shrinkage matters. For $t>\mu$, the dimensionless ratio
$\kappa(t-\mu)/(n\varepsilon)>1$ means that shrinkage raises the threshold
above the harmed cell's expected error count. This is a finite-sample
diagnostic, not a minimax phase transition. The integer cutoff may remain
unchanged after a small increase in the real threshold.

An exact example illustrates this separation. Take five equally sized
cells with $n=10$, four with rate $0.1$ and one with rate $0.35$, and fix
$\mu=0.1$, $\kappa=90$, and $\Delta=0.1$. The exceptional cell's complement
has rate $r=0.1$, so its oracle threshold is $t=0.2$. Macro mean squared error (MSE) is
$0.01175$ for raw proportions and $0.0102425$ for shrinkage. The raw
screen selects this cell with probability
$\Pr_{0.35}\{X\ge3\}=0.7383926086$, while the shrunk screen would require
$X>11$ and has selection probability zero. These numbers follow from
finite-binomial calculations rather than fitted experiment outcomes. The
prior is deliberately strong; whether fitted methods reach a comparable
regime is an empirical question.

More generally, each cell with $p=\mu$ has strictly smaller risk under
this shrinker. Sufficiently many such cells offset the finite risk cost of
one exceptional cell, even when its screen collapses. The calculation
holds conditionally for hyperparameters learned from independent data.
Treating hyperparameters fitted to the same data as fixed would not justify
Eq.~\eqref{eq:theory-risk}. Full proofs appear in
Appendix~\ref{app:theory-proofs}.

\section{Confirmation and allocation under a common budget}
\label{sec:confirmation}

\subsection{Unknown-reference confirmation}

Our primary confirmation engine combines a Hoeffding upper bound for the
reference with an exact binomial target tail; we refer to it as HB. The
ingredients are established concentration and nuisance-parameter
methods~\citep{hoeffding1963probability,berger1994pvalues}.

\begin{proposition}[HB p-value with a heterogeneous reference]
\label{prop:theory-hb}
Let $X\sim\mathrm{Bin}(n,p)$, $n\ge1$. Let the reference observations
$B_1,\ldots,B_m$, $m\ge1$, be independent Bernoulli variables with
possibly unequal means $q_1,\ldots,q_m$. Define
$q=m^{-1}\sum_jq_j$ and $\widehat q=m^{-1}\sum_j B_j$.
Fix $0<\eta<1$ before observing confirmation outcomes, and set
\begin{equation}
 \begin{aligned}
 U_q&=\min\left\{1,\widehat q+
          \sqrt{\frac{\log(1/\eta)}{2m}}\right\},\\
 \theta_U&=\min\{1,U_q+\Delta\}.
 \end{aligned}
 \label{eq:theory-upper-reference}
\end{equation}
\begin{equation}
 \begin{aligned}
 P_{\mathrm{HB}}&=\min\{1,\eta+T_{\theta_U}(X)\},\\
 T_\theta(x)&=\Pr\{\mathrm{Bin}(n,\theta)\ge x\}.
 \end{aligned}
 \label{eq:theory-hb}
\end{equation}
Under $p-q\le\Delta$, $\Pr(P_{\mathrm{HB}}\le u)\le u$ for all
$u\in[0,1]$. Setting $P_{\mathrm{HB}}=1$ when either denominator is zero
preserves validity.
\end{proposition}

The binomial target assumption is essential: an atomic cell has a common
rate, whereas an aggregate of heterogeneous cells need not. The proof
allows a heterogeneous complement and does not replace its law by a
binomial approximation. For a fixed standardized reference
$q=\sum_h\nu_hp_h$, with fixed $\nu_h\ge0$ and $\sum_h\nu_h=1$,
samples from supported cells instead give
$\widehat q=\sum_h\nu_h\widehat p_h$ and upper radius
$\{\frac12\log(1/\eta)\sum_h\nu_h^2/n_h\}^{1/2}$.
Unsupported cells with positive reference weight require a worst-case
contribution or an insufficient-information outcome.

For heterogeneous targets, a separate valid alternative uses independent
Bernoulli observations on both sides. With empirical difference $d$ and
$V=1/n+1/m$, the difference-Hoeffding p-value is
\begin{equation}
 \begin{aligned}
 P_{\mathrm H}&=\exp\{-2(d-\Delta)_+^2/V\},\\
 (x)_+&=\max\{x,0\}.
 \end{aligned}
 \label{eq:theory-hoeffding}
\end{equation}
It tests the difference of the two average means. We use its conservative
boundary for planning below; it does not supply a power formula for HB.

\subsection{Independent discovery weights and weighted Holm}

Discovery chooses nonnegative scores $s_g(D)$ without using confirmation
outcomes. For a fixed floor $0<\rho\le1$, define
\begin{equation}
 w_g=\frac{\rho}{G}+(1-\rho)\frac{s_g}{\sum_hs_h},
 \label{eq:theory-weights}
\end{equation}
using $w_g=1/G$ when all scores vanish. Thus $w_g\ge\rho/G$ and
$\sum_gw_g=1$. Independent weighting and sample splitting are established
multiple-testing tools~\citep{kang2009weighted,rubin2006splitting}; here
their role is to allocate a fixed audit budget.

\begin{proposition}[Conditional strong FWER control]
\label{prop:theory-holm}
Fix $0<\alpha<1$. Suppose the confirmation p-values for true nulls are super-uniform conditional on
discovery and the legitimate confirmation design. Order hypotheses by
$P_g/w_g$, with a fixed rule for ties, and reject sequentially while
\begin{equation}
 P_{(j)}\le\frac{\alpha w_{(j)}}{\sum_{k=j}^G w_{(k)}},
 \label{eq:theory-holm}
\end{equation}
stopping at the first failure. Then FWER is at most $\alpha$, allowing
arbitrary dependence among confirmation p-values. The rejection set
contains every weighted-Bonferroni rejection $P_g\le\alpha w_g$.
\end{proposition}

The proof is the usual first-false-rejection argument for step-down
testing~\citep{holm1979simple}. Misspecified discovery scores affect
allocation, while conditional validity provides error control. The floor
ensures a positive minimum testing budget; it gives neither equal power
nor equal information.

The five score options are equal scores, $\sqrt{n_g^D}$, raw excess
$(\widehat\delta_g^D-\Delta)_+$, empirical Bayes excess
$(\widetilde\delta_g^D-\Delta)_+$, and a conflict score
\begin{equation}
 \begin{aligned}
 s_g^{\mathrm{conflict}}&=(\widehat\delta_g^D-\Delta)_+\\
 &\quad+(\widehat\delta_g^D-\widetilde\delta_g^D)_+.
 \end{aligned}
 \label{eq:theory-conflict}
\end{equation}
Raw and fitted differences use their respective count-weighted
complements. Scores are zero for cells without discovery observations.
The conflict score is a baseline option for testing whether disagreement
helps allocation; a consistent additional benefit has not been demonstrated, and the score is not
presented as a new testing principle.

The primary comparator uses all audit observations with HB and ordinary
Holm, corresponding to equal weights. The split methods use the same total
count budget, estimand, and p-value engine. To isolate floor effects, all
main comparisons and floor ablations fix
$\eta=\alpha\rho_{\min}/(20G)$ with $\rho_{\min}=0.05$, rather than changing
the HB engine with $\rho$. For planned floors $\rho\ge\rho_{\min}$,
$\eta\le\alpha w_g/20$. Score definitions, floors, and nuisance budgets
are frozen before confirmation.

\subsection{The allocation advantage must offset splitting}

For supported confirmation counts, let $V_g=1/n_g^C+1/m_g^C$.
For Eq.~\eqref{eq:theory-hoeffding}, define the sufficient boundary
$b_g(w)=\{(V_g/2)\log[1/(\alpha w)]\}^{1/2}$. If
$\delta_g=\Delta+\varepsilon_g$ and $\varepsilon_g>b_g(w_g)$,
weighted Holm has conditional power at least
\begin{equation}
 1-\exp\{-2[\varepsilon_g-b_g(w_g)]^2/V_g\}.
 \label{eq:theory-power-envelope}
\end{equation}
Below that boundary the stated lower bound is zero, which says nothing
about whether actual power is zero. Replacing $w_g$ by $\rho/G$ gives a
score-independent version.

\begin{corollary}[Comparison of Hoeffding sufficient boundaries]
\label{cor:theory-split}
Suppose confirmation retains exactly a fraction $c$ of each relevant count,
so $V_C=V_A/c$, where $A$ denotes all audit data. The split
weighted-Bonferroni sufficient boundary is strictly smaller than the
full-data equal-weight Bonferroni boundary if and only if
\begin{equation}
 w_g>\alpha^{c-1}G^{-c}.
 \label{eq:theory-split-cost}
\end{equation}
\end{corollary}

At $\alpha=0.05$ and $c=0.5$, this requires $w_g>0.5590$ for $G=64$,
and $w_g>1.1180$ for $G=16$. The latter is impossible. These are
comparisons of conservative sufficient boundaries, not an ordering of
actual powers, and do not compare split and full-data Holm directly.
Full-data Holm can improve upon its Bonferroni counterpart. Consequently,
the FWER guarantee provides no general power advantage for splitting;
that question requires the paired empirical comparisons.
